\chapter{Методы используемые в исследовании спектров молекул}\label{ch:methods}

\section{Оператор Гамильтона молекулы}\label{sec:methods/hamilton}

\subsection{Роль гамильтониана в квантово-механическом описании молекул}

Квантово-механическое описание молекулярных систем основано на использовании
оператора Гамильтона, определяющего полную энергию системы.
В стационарной квантовой механике состояние молекулы описывается
волновой функцией $\Psi$, удовлетворяющей уравнению Шрёдингера

\begin{equation}
\hat{H}\Psi = E\Psi ,
\end{equation}

где $\hat{H}$ — оператор Гамильтона, а $E$ — собственное значение энергии.

Собственные значения гамильтониана соответствуют энергетическим уровням
молекулы, а переходы между ними приводят к появлению линий в спектрах
поглощения и излучения. Таким образом, анализ молекулярных спектров
фактически сводится к исследованию структуры собственных значений
гамильтониана.

Формирование гамильтонова подхода в молекулярной физике связано
с развитием квантовой механики и спектроскопии в первой половине XX века.
Фундаментальные основы теории колебательно-вращательных спектров
были сформулированы в классических работах
\cite{nielsen,Amat1971,Makushkin1984}.
Эти исследования показали, что даже относительно простые молекулярные
системы обладают сложной энергетической структурой,
которая определяется взаимосвязью вращательного и колебательного движения.

\subsection{Полный гамильтониан молекулы}

В нерелятивистском приближении гамильтониан молекулы,
состоящей из $N$ электронов и $M$ атомных ядер, можно записать как

\begin{equation}
\hat{H} =
-\sum_{i=1}^{N}\frac{\hbar^2}{2m_e}\nabla_i^2
-\sum_{A=1}^{M}\frac{\hbar^2}{2M_A}\nabla_A^2
+\sum_{i<j}\frac{e^2}{r_{ij}}
-\sum_{i,A}\frac{Z_A e^2}{r_{iA}}
+\sum_{A<B}\frac{Z_A Z_B e^2}{R_{AB}} .
\end{equation}

Здесь первые два члена описывают кинетическую энергию электронов
и атомных ядер, а остальные члены учитывают кулоновские взаимодействия
между всеми частицами системы.

Прямое решение уравнения Шрёдингера с таким гамильтонианом
для многоатомных молекул представляет собой чрезвычайно сложную задачу.
Поэтому в практических расчётах используются приближённые методы,
основой которых является приближение Борна–Оппенгеймера.
В этом приближении движение электронов и ядер рассматривается раздельно,
что позволяет сначала решить электронную задачу,
а затем исследовать движение ядер в эффективном потенциале.

После отделения движения центра масс остаётся задача описания
внутреннего движения ядер, включающего вращательные
и колебательные степени свободы.

\subsection{Колебательно-вращательный гамильтониан}

Внутреннее движение ядер молекулы характеризуется
$3N-6$ колебательными координатами
(для нелинейной молекулы) и тремя вращательными степенями свободы.
Соответствующий гамильтониан можно представить в виде

\begin{equation}
\hat{H} =
\hat{H}_{vib} +
\hat{H}_{rot} +
\hat{H}_{vr},
\end{equation}

где $\hat{H}_{vib}$ описывает колебательное движение,
$\hat{H}_{rot}$ — вращательное движение,
а $\hat{H}_{vr}$ учитывает взаимодействие между ними.

В гармоническом приближении колебательная часть имеет вид

\begin{equation}
\hat{H}_{vib} =
\frac{1}{2}\sum_k
\left(P_k^2 + \omega_k^2 Q_k^2 \right),
\end{equation}

где $Q_k$ — нормальные координаты,
$P_k$ — сопряжённые импульсы,
а $\omega_k$ — частоты нормальных колебаний.

Вращательная часть гамильтониана для нелинейной молекулы
может быть записана как

\begin{equation}
\hat{H}_{rot} =
A\hat{J}_a^2 +
B\hat{J}_b^2 +
C\hat{J}_c^2 ,
\end{equation}

где $A$, $B$ и $C$ — вращательные постоянные,
связанные с главными моментами инерции молекулы.

\subsection{Операторы углового момента}

Вращательное движение молекулы описывается оператором
полного углового момента $\hat{\mathbf{J}}$.
Его компоненты $\hat{J}_a$, $\hat{J}_b$, $\hat{J}_c$
в молекулярной системе координат удовлетворяют
коммутационным соотношениям

\begin{equation}
[\hat{J}_i,\hat{J}_j] = i\hbar \varepsilon_{ijk}\hat{J}_k .
\end{equation}

Собственные функции операторов $\hat{J}^2$ и $\hat{J}_z$
определяются квантовыми числами $J$ и $M$:

\begin{align}
\hat{J}^2 |J,M\rangle &= \hbar^2 J(J+1)|J,M\rangle, \\
\hat{J}_z |J,M\rangle &= \hbar M |J,M\rangle .
\end{align}

Квантовое число $J$ характеризует полный момент вращения молекулы,
а $M$ определяет его проекцию на выбранную ось.

Для симметричных молекул дополнительно вводится
квантовое число $K$, определяющее проекцию
углового момента на ось симметрии молекулы.

\subsection{Симметрия волновых функций}

Волновая функция молекулы может быть представлена
в виде произведения

\begin{equation}
\Psi = \Psi_{el}\Psi_{vib}\Psi_{rot}.
\end{equation}

Симметрия каждой из этих функций определяется
точечной группой симметрии молекулы.
Использование свойств симметрии позволяет существенно
уменьшить размерность квантово-механической задачи,
поскольку гамильтониан коммутирует с операциями
группы симметрии.

Симметрия состояний играет важную роль
в формулировке правил отбора для спектроскопических переходов
и при классификации энергетических уровней молекулы.

\subsection{Вывод гамильтониана Уотсона}

При описании движения ядер в молекуле
необходимо учитывать как вращение,
так и колебания атомов.
Полное разделение этих движений невозможно,
однако их связь можно минимизировать
путём выбора специальной системы координат,
известной как система Экарта \cite{Eckart1935,BunkerJensen1998}.
В такой системе координат оси связаны с молекулой
и изменяются вместе с её колебаниями.

После применения условий Экарта
гамильтониан движения ядер можно преобразовать
в более удобную форму.
В 1968 году Дж.~Уотсон предложил преобразование,
позволяющее существенно упростить
колебательно-вращательный гамильтониан молекулы
\cite{Watson1968}.

В результате гамильтониан движения ядер
можно записать в виде

\begin{equation}
\hat{H} =
-\frac{\hbar^2}{2}\sum_{s=1}^{3N-6}
\frac{\partial^2}{\partial q_s^2}
+
\frac{1}{2}
\sum_{\alpha\beta}
\mu_{\alpha\beta}
(\hat{J}_\alpha - \hat{\pi}_\alpha)
(\hat{J}_\beta - \hat{\pi}_\beta)
+V(q_1,\ldots,q_{3N-6}),
\end{equation}

где

\begin{itemize}
\item $q_s$ — нормальные колебательные координаты,
\item $\hat{J}_\alpha$ — компоненты оператора углового момента,
\item $\hat{\pi}_\alpha$ — операторы кориолисового взаимодействия,
\item $\mu_{\alpha\beta}$ — элементы обратного тензора инерции,
\item $V$ — потенциальная энергия.
\end{itemize}

Полученный гамильтониан известен как гамильтониан Уотсона
и широко используется при анализе
колебательно-вращательных спектров молекул
\cite{Watson1968,Watson1977}.

\subsection{Переход к матричному представлению гамильтониана}

Для практического вычисления энергетических уровней
гамильтониан представляется в виде матрицы
в выбранном базисе колебательно-вращательных функций

\begin{equation}
H_{ij} =
\langle \Phi_i | \hat{H} | \Phi_j \rangle .
\end{equation}

Диагонализация этой матрицы позволяет получить
собственные значения энергии и соответствующие
волновые функции молекулы.

Построение эффективных базисных функций
и вычисление матричных элементов гамильтониана
рассматриваются в следующем разделе,
где будет подробно описано
матричное представление гамильтониана
и методы его использования
при анализе колебательно-вращательных спектров.

\section{Приближение Борна-Оппенгеймера}

\subsection{Физический смысл и предпосылки приближения}

При квантово-механическом описании молекулы необходимо одновременно
учитывать движение электронов и атомных ядер. Полная волновая функция
молекулярной системы зависит от координат всех частиц, а соответствующее
уравнение Шрёдингера содержит как электронную, так и ядерную кинетическую
энергию, а также все кулоновские взаимодействия. Для многоатомных
молекул такая задача в исходной постановке обладает чрезвычайно высокой
размерностью, вследствие чего её прямое решение оказывается практически
невыполнимым даже для сравнительно простых систем.

Ключевое упрощение связано с тем, что атомные ядра существенно тяжелее
электронов. Именно это различие масштабов масс приводит к различию
характерных времён движения соответствующих подсистем. Электронное
облако успевает перестроиться значительно быстрее, чем заметно изменяется
положение ядер. Иными словами, для электронной подсистемы конфигурацию
ядер в первом приближении можно считать фиксированной, тогда как для
ядерного движения электронная энергия играет роль эффективного потенциала.
Такой подход и составляет физическую основу приближения
Борна--Оппенгеймера.

Исторически данное приближение стало одним из важнейших результатов
раннего этапа развития квантовой механики. Оно позволило перейти
от формально точного, но практически непригодного полного уравнения
Шрёдингера к иерархической схеме описания молекулы, в которой
электронная, колебательная и вращательная задачи рассматриваются
последовательно. Благодаря этому приближение Борна--Оппенгеймера
фактически стало фундаментом современной квантовой химии,
молекулярной физики и спектроскопии высокого разрешения.

\subsection{Разделение полной молекулярной задачи}

Пусть $\mathbf{r}$ обозначает совокупность электронных координат,
а $\mathbf{R}$ --- совокупность координат атомных ядер.
Тогда полная волновая функция молекулы имеет вид
$\Psi(\mathbf{r},\mathbf{R})$ и удовлетворяет стационарному уравнению

\begin{equation}
\hat{H}\Psi(\mathbf{r},\mathbf{R})=E\Psi(\mathbf{r},\mathbf{R}).
\end{equation}

В рамках приближения Борна--Оппенгеймера предполагается, что полную
волновую функцию можно представить в виде произведения электронной
и ядерной частей:

\begin{equation}
\Psi(\mathbf{r},\mathbf{R})
=
\psi_{el}(\mathbf{r};\mathbf{R})\chi(\mathbf{R}).
\end{equation}

Здесь электронная функция $\psi_{el}(\mathbf{r};\mathbf{R})$
зависит от координат ядер параметрически:
при решении электронной задачи положения ядер фиксируются,
а их координаты входят в уравнение лишь как внешние параметры.
Ядерная функция $\chi(\mathbf{R})$ описывает уже более медленное
движение самих ядер.

Подстановка такого разложения в полное уравнение Шрёдингера приводит
к естественному разделению задачи на два этапа.
Сначала решается электронное уравнение при фиксированной конфигурации ядер:

\begin{equation}
\hat{H}_{el}(\mathbf{r};\mathbf{R})
\psi_{el}(\mathbf{r};\mathbf{R})
=
E_{el}(\mathbf{R})\psi_{el}(\mathbf{r};\mathbf{R}),
\end{equation}

где $\hat{H}_{el}$ включает кинетическую энергию электронов,
электрон-электронное взаимодействие, электрон-ядерное взаимодействие
и межъядерное отталкивание как параметрическую величину.
Для каждой фиксированной конфигурации ядер получается собственное
значение $E_{el}(\mathbf{R})$, которое зависит от геометрии молекулы.

Далее найденная электронная энергия используется в уравнении
для ядерного движения:

\begin{equation}
\left[
-\sum_A\frac{\hbar^2}{2M_A}\nabla_A^2
+E_{el}(\mathbf{R})
\right]\chi(\mathbf{R})
=
E\chi(\mathbf{R}),
\end{equation}

если не учитывать более тонкие поправки, возникающие из-за
неполного разделения электронного и ядерного движения.
Тем самым исходная многочастичная задача заменяется последовательным
решением электронной и ядерной задач, что делает возможным
практический анализ молекулярных состояний.

\subsection{Потенциальные поверхности энергии и их роль}

Функция $E_{el}(\mathbf{R})$, получаемая как результат решения
электронного уравнения, называется потенциальной поверхностью энергии.
Именно она определяет, в каком эффективном поле движутся ядра,
а значит, задаёт геометрию молекулы, её колебательные свойства
и значительную часть наблюдаемой спектроскопической структуры.

Минимум потенциальной поверхности соответствует равновесной конфигурации
молекулы. В окрестности этого минимума движение ядер может быть
рассмотрено как малые отклонения от равновесия, что приводит
к введению нормальных координат и гармоническому приближению
для колебаний. Кривизна поверхности вблизи минимума определяет
гармонические частоты, а члены более высокого порядка описывают
ангармоничность и взаимодействие колебательных мод.

Если в системе существует несколько электронных состояний,
то каждому из них соответствует своя потенциальная поверхность.
В большинстве спектроскопических задач высокого разрешения
рассматривается движение ядер на одной выбранной электронной поверхности,
обычно соответствующей основному электронному состоянию.
Именно в этом случае приближение Борна--Оппенгеймера проявляет
свою наибольшую эффективность.

Следует подчеркнуть, что в молекулярной спектроскопии
потенциальная поверхность энергии является не просто промежуточной
математической конструкцией, а фактически центральным объектом
описания внутренней динамики молекулы. Через неё определяются
равновесные структурные параметры, колебательные уровни,
резонансные взаимодействия и, в конечном счёте, положение
наблюдаемых спектральных линий.

\subsection{Адиабатическое приближение и структура молекулярной энергии}

При использовании приближения Борна--Оппенгеймера молекулярную энергию
естественно рассматривать как величину, распадающуюся на несколько
компонент, соответствующих различным типам движения. В первом
приближении полная энергия молекулы может быть представлена как сумма
электронного, колебательного и вращательного вкладов. Такое разложение
не является строго точным, однако именно оно лежит в основе всей
классификации молекулярных уровней и спектров.

Сначала формируется электронная структура молекулы, определяемая
решением электронной задачи. Затем на фиксированной электронной
поверхности возникают колебательные уровни, связанные с движением
ядер относительно равновесной конфигурации. Наконец, каждый
колебательный уровень расщепляется на вращательные подуровни.
Именно эта иерархия масштабов энергии делает возможным
последовательное построение молекулярного гамильтониана
и поэтапный анализ спектров.

Такой подход особенно важен для колебательно-вращательной
спектроскопии. Наблюдаемые в инфракрасных спектрах переходы,
как правило, происходят в пределах одного электронного состояния
и отражают изменение колебательных и вращательных квантовых чисел.
Следовательно, сам факт возможности анализировать эти переходы
независимо от электронной структуры в значительной степени основан
именно на приближении Борна--Оппенгеймера.

\subsection{Математические основания и малый параметр}

С формальной точки зрения приближение Борна--Оппенгеймера опирается
на существование малого параметра, связанного с отношением массы
электрона к характерной массе ядра. Благодаря этому полное уравнение
Шрёдингера может быть исследовано асимптотически, а электронное
и ядерное движения оказываются разделимыми в нулевом порядке
по указанному малому параметру.

Именно в этом смысле приближение Борна--Оппенгеймера не является
лишь качественным рассуждением о ``быстрых'' электронах и
``медленных'' ядрах, а представляет собой вполне определённую
схему последовательного приближения. В первом порядке получаются
разделённые электронная и ядерная задачи. В более высоких порядках
возникают поправки, связанные с тем, что электронная волновая функция
всё же зависит от координат ядер не только параметрически, но и
через производные по этим координатам.

Эти дополнительные члены обычно малы, однако в ряде случаев
они могут оказывать заметное влияние на спектр. Поэтому
в прикладной спектроскопии важно не только использовать
приближение Борна--Оппенгеймера как исходный уровень описания,
но и понимать границы его применимости.

\subsection{Нарушение приближения и неадиабатические эффекты}

Хотя приближение Борна--Оппенгеймера чрезвычайно эффективно,
оно не является строго точным. Основная причина заключается в том,
что электронная функция зависит от координат ядер, а потому
при действии ядерного кинетического оператора появляются дополнительные
члены, связывающие разные электронные состояния. Эти члены и образуют
так называемые неадиабатические связи.

В обычных условиях, когда электронные состояния хорошо разделены по энергии,
такие связи малы, и ими можно пренебречь. Однако ситуация меняется,
если потенциальные поверхности сближаются или пересекаются.
В этом случае разделение электронного и ядерного движения ухудшается,
а простая одноповерхностная картина становится недостаточной.
Подобные эффекты особенно важны в фотохимии, в теории элементарных
реакций, а также при анализе некоторых тонких спектроскопических
особенностей высоковозбуждённых состояний.

Для задач, связанных с высокоразрешённой инфракрасной спектроскопией
устойчивых молекул в основном электронном состоянии, нарушение
приближения Борна--Оппенгеймера обычно не является определяющим.
Тем не менее именно понимание существования таких поправок
объясняет, почему даже очень точные эффективные гамильтонианы
следует рассматривать как приближённые модели, а не как абсолютно
полное описание молекулярной динамики.

\subsection{Значение приближения Борна--Оппенгеймера для молекулярной спектроскопии}

Для молекулярной спектроскопии приближение Борна--Оппенгеймера
имеет принципиальное значение. Именно оно позволяет перейти
от общей многочастичной задачи к последовательному описанию
электронных, колебательных и вращательных состояний, что делает
возможной интерпретацию экспериментальных спектров в терминах
квантовых чисел и эффективных спектроскопических параметров.

В частности, после фиксации электронной поверхности энергии
задача сводится к исследованию движения ядер. Далее вводятся
нормальные координаты, строится колебательный гамильтониан,
а затем учитывается вращательное движение и его связь с колебаниями.
Именно в этой логике формируются как строгие квантово-механические,
так и полуэмпирические методы анализа спектров высокого разрешения.

Таким образом, приближение Борна--Оппенгеймера выступает не просто
как техническое упрощение исходного уравнения Шрёдингера, а как
фундаментальный принцип организации молекулярной теории.
Оно связывает электронную структуру молекулы с её ядерной динамикой
и тем самым создаёт теоретическую основу для последующего рассмотрения
колебательных состояний, колебательно-вращательного взаимодействия
и построения эффективных гамильтонианов, используемых
в практической спектроскопии.

В дальнейшем, опираясь на данное приближение, можно перейти
к описанию колебательного движения ядер в окрестности равновесной
конфигурации молекулы. Именно этот этап приводит к введению
нормальных координат, колебательных волновых функций и,
в конечном счёте, к матричному представлению молекулярного
гамильтониана, необходимому для анализа экспериментальных спектров.

\section{Теория изотопозамещения}

В задачах молекулярной спектроскопии под изотопозамещением понимают замену одного
или нескольких ядер молекулы их изотопами при сохранении химического состава и
электронной структуры. В рамках приближения Борна--Оппенгеймера такая замена
не изменяет электронную потенциальную поверхность как функцию декартовых
координат ядер, поскольку зарядовое распределение ядер остается тем же. Однако
массы ядер входят в оператор кинетической энергии, поэтому меняются моменты
инерции, нормальные координаты, частоты нормальных колебаний, кориолисовы
параметры и, как следствие, положения и интенсивности колебательно-вращательных
линий. Именно это обстоятельство делает изотопологи одновременно близкими к
``материнской'' молекуле и достаточно различными, чтобы их спектры давали
дополнительную информацию о внутримолекулярной динамике
\cite{Wilson1936,WilsonDeciusCross1955,Papousek1982,BunkerJensen1998}.

Для основной модификации молекулы ядерная часть гамильтониана в декартовых
координатах может быть записана в виде
\begin{equation}
\hat{H}=\sum_{N,\alpha}\frac{\hat{P}_{N\alpha}^{2}}{2m_N}+V(x),
\end{equation}
где $m_N$ --- масса $N$-го ядра, $\hat{P}_{N\alpha}$ --- соответствующая компонента
оператора импульса, $x$ --- совокупность декартовых координат ядер, а $V(x)$ ---
потенциальная энергия молекулы. Для изотополога аналогичный оператор имеет вид
\begin{equation}
\hat{H}'=\sum_{N,\alpha}\frac{\hat{P}_{N\alpha}^{2}}{2m'_N}+V(x)
       =\hat{H}+\Delta \hat{T},
\end{equation}
где
\begin{equation}
\Delta \hat{T}=\sum_{N,\alpha}
\left(\frac{1}{2m'_N}-\frac{1}{2m_N}\right)\hat{P}_{N\alpha}^{2} .
\end{equation}
Тем самым изотопический эффект в первом приближении сводится к изменению
кинетической части гамильтониана. Потенциальная функция в декартовом
представлении считается общей для всех изотопологов, а различия возникают при
переходе к масс-зависимым координатам, в которых решается колебательно-вращательная задача \cite{byk}.

Практическая теория изотопозамещения строится на сравнении двух систем
координат: системы, связанной с основной модификацией молекулы, и системы,
связанной с изотопологом. При замене масс изменяется положение центра масс и,
в общем случае, ориентация главных осей инерции. Поэтому даже если равновесная
геометрия молекулы в декартовом пространстве считается неизменной, координаты
ядер в молекулярно-фиксированной системе изотополога оказываются другими.
Это особенно важно для несимметричного изотопозамещения: например, переход от
$\mathrm{H_2S}$ или $\mathrm{D_2S}$ к $\mathrm{HDS}$ сопровождается понижением симметрии
и поворотом молекулярной системы координат, что непосредственно отражается на
структуре гамильтониана и на виде оператора дипольного момента
\cite{hds_sydow2019}.

Связь между нормальными координатами основной модификации и изотополога можно
представить в виде разложения
\begin{equation}
Q_\lambda=\sum_{\mu}a_{\mu\lambda}Q'_\mu+
\sum_{\mu\nu}a^{\lambda}_{\mu\nu}Q'_\mu Q'_\nu+\ldots,
\end{equation}
где коэффициенты $a_{\mu\lambda}$ и $a^{\lambda}_{\mu\nu}$ определяются массами ядер,
формами нормальных колебаний и выбором молекулярно-фиксированных осей.
В гармоническом приближении эти коэффициенты позволяют связать частоты
изотополога с частотами основной модификации. В более высоких порядках они
задают преобразование кубических, квартетных и последующих членов потенциальной
функции. Поэтому в нормальных координатах параметры разложения потенциальной
энергии становятся изотопически зависимыми, хотя сама поверхность $V(x)$ в
приближении Борна--Оппенгеймера является общей.

Если потенциальная энергия в окрестности равновесной конфигурации представлена
в виде
\begin{equation}
V=\frac{1}{2}\sum_i \omega_i^2 Q_i^2+
\frac{1}{6}\sum_{ijk}K_{ijk}Q_iQ_jQ_k+
\frac{1}{24}\sum_{ijkl}K_{ijkl}Q_iQ_jQ_kQ_l+\ldots,
\end{equation}
то замена $Q_i\rightarrow Q'_i$ приводит к новому набору величин
$\omega'_i$, $K'_{ijk}$, $K'_{ijkl}$ и т.д. На уровне гармонического приближения
это эквивалентно изменению матрицы кинетической энергии в задаче Вильсона
$GF$; силовая матрица в внутренних или декартовых координатах при этом
сохраняется, а матрица $G$ изменяется вследствие замены масс. Поэтому
характерные изотопические сдвиги частот могут быть оценены уже из векового
уравнения для нормальных колебаний. Классические соотношения для вибрационных
изотопических эффектов показывают, что отдельные частоты зависят от силового
поля, тогда как некоторые комбинации частот, масс и моментов инерции обладают
более универсальным характером \cite{SalantRosenthal1932,WilsonDeciusCross1955}.

Изменение моментов инерции приводит также к изменению вращательных постоянных
\begin{equation}
A'=\frac{h}{8\pi^2 c I'_a},\qquad
B'=\frac{h}{8\pi^2 c I'_b},\qquad
C'=\frac{h}{8\pi^2 c I'_c}.
\end{equation}
Для тяжелого изотопозамещения, особенно при замене водорода дейтерием, эффект
оказывается значительным: возрастают моменты инерции и, как правило,
уменьшаются вращательные постоянные. Поэтому линии изотополога смещаются не
только из-за изменения колебательной энергии, но и вследствие перестройки
вращательной структуры. В асимметричных волчках это дополнительно влияет на
параметр асимметрии и на степень смешивания базисных функций с различными
значениями $K_a$ и $K_c$.

Не менее важным следствием изотопозамещения является изменение симметрии
молекулы. Если замена проводится в эквивалентной позиции и не нарушает точечную
группу, то правила отбора и типы полос сохраняются, хотя численные значения
спектроскопических параметров меняются. Если же замена понижает симметрию,
часть вырождений снимается, ранее эквивалентные колебания становятся
неэквивалентными, а некоторые переходы, запрещенные или слабые в материнской
молекуле, могут проявляться в спектре изотополога. Для молекул типа
$\mathrm{XY_4}$ это особенно заметно при переходе от сферического волчка к
изотопологам более низкой симметрии; для молекул типа $\mathrm{XY_2}$ примером
служит переход $\mathrm{H_2S}\rightarrow\mathrm{HDS}$, при котором изменяются и
симметрия, и ориентация главных осей инерции \cite{d2s_sydow2019,hds_sydow2019,4300}.

Теория изотопозамещения используется не только для предсказания положений
линий, но и для анализа их интенсивностей. В колебательно-вращательной
спектроскопии интенсивности определяются матричными элементами оператора
дипольного момента. В декартовом представлении поверхность дипольного момента,
как и потенциальная поверхность, в приближении Борна--Оппенгеймера является
общей для изотопологов. Однако ее разложение по нормальным координатам и
проекции на молекулярно-фиксированные оси изменяются. Если для материнской
молекулы
\begin{equation}
\mu_\alpha=\mu^e_\alpha+\sum_i \mu^i_\alpha Q_i+
\frac{1}{2}\sum_{ij}\mu^{ij}_\alpha Q_iQ_j+\ldots,
\end{equation}
то для изотополога необходимо учитывать как преобразование нормальных координат,
так и возможный поворот осей:
\begin{equation}
\mu'_\alpha=\sum_\beta K_{\alpha\beta}(Q')
\left(\mu^e_\beta+\sum_i \mu^i_\beta Q_i(Q')+
\frac{1}{2}\sum_{ij}\mu^{ij}_\beta Q_i(Q')Q_j(Q')+\ldots\right).
\end{equation}
Отсюда следуют изотопические соотношения между параметрами эффективного
дипольного момента. Именно эти соотношения позволяют использовать хорошо
изученную основную модификацию как источник исходной информации для расчета
интенсивностей менее изученных изотопологов \cite{Ulenikov2019a,Ulenikov2020a}.

В настоящей работе этот подход имеет принципиальное значение. При исследовании
изотопологов сероводорода и тетрагидридов элементов IV группы экспериментальная
смесь может содержать несколько изотопических модификаций, а их парциальные
давления не всегда могут быть определены напрямую с требуемой точностью. Тогда
используется обратная схема: параметры эффективного дипольного момента
изотополога оцениваются из параметров материнской молекулы, после чего
рассчитанные интенсивности сравниваются с экспериментальными. Масштабный
множитель между рассчитанными и измеренными интенсивностями дает возможность
оценить относительное содержание изотополога в газовой смеси. Такой прием был
реализован в работах по $\mathrm{D_2S}$ и $\mathrm{HDS}$, где изотопические
соотношения использовались при анализе полосы $\nu_2$ и при определении
параметров эффективного дипольного момента \cite{d2s_sydow2019,hds_sydow2019}.

Для молекул высокой симметрии роль изотопозамещения проявляется несколько
иначе. В изотопологах типа $^{M}\mathrm{GeH_4}$ или $^{M}\mathrm{SiH_4}$, где
заменяется центральное ядро, точечная группа $T_d$ сохраняется, но меняются
массы, частоты и параметры эффективного гамильтониана. Это дает возможность
согласованно анализировать несколько изотопологов в рамках одной потенциальной
модели и уточнять параметры колебательной структуры. Если же изотопозамещение
затрагивает периферийные атомы, симметрия обычно понижается, что требует
дополнительного учета расщеплений, новых правил отбора и преобразований
базисных функций. В обоих случаях информация об изотопологах служит
дополнительным ограничением для потенциальной поверхности и для эффективных
моделей, используемых при расчете спектров \cite{872,4300}.

Таким образом, теория изотопозамещения связывает спектроскопические параметры
разных изотопических модификаций одной молекулы через общность потенциальной и
дипольной поверхностей и через явную зависимость ядерной динамики от масс.
Она позволяет переносить информацию от хорошо изученной материнской молекулы к
менее изученным изотопологам, оценивать частоты, вращательные постоянные,
параметры эффективного гамильтониана и эффективного дипольного момента, а также
контролировать состав изотопической газовой смеси по спектральным интенсивностям.
По этой причине теория изотопозамещения является естественным связующим звеном
между общими положениями колебательно-вращательной спектроскопии и практическими
методами анализа спектров, которые используются в последующих главах.

\section{Неприводимые тензорные операторы}\label{sec:methods/ito}

\subsection{Общее определение и связь с симметрией вращений}

В задачах молекулярной спектроскопии операторный формализм используется не
только для записи гамильтониана и оператора дипольного момента, но и для
выявления их симметрийных свойств. При прямой записи через декартовы
компоненты операторы быстро становятся громоздкими, особенно если необходимо
учитывать вырождение колебаний, кориолисовы взаимодействия, центробежные
искажения и резонансные связи между состояниями. Более удобной оказывается
запись в терминах неприводимых тензорных операторов, то есть таких наборов
операторных компонент, которые при вращениях преобразуются как одно
неприводимое представление группы вращений. Такой подход восходит к теории
углового момента и неприводимых тензорных наборов Фано--Ракаха и Вигнера
\cite{Fano1959,Wigner1965,Varshalovich1988} и в дальнейшем стал одним из
стандартных инструментов теории колебательно-вращательных спектров
многоатомных молекул.

Пусть $\hat{T}^{(k)}_q$ обозначает компоненту сферического тензорного
оператора ранга $k$, где $q=-k,-k+1,\ldots,k$. Набор
$\{\hat{T}^{(k)}_q\}$ называется неприводимым, если при произвольном
повороте системы координат $\Omega$ его компоненты преобразуются по закону

\begin{equation}
\hat{R}(\Omega)\hat{T}^{(k)}_q\hat{R}^{-1}(\Omega)=
\sum_{q'=-k}^{k}D^{(k)}_{q'q}(\Omega)\hat{T}^{(k)}_{q'},
\end{equation}

где $\hat{R}(\Omega)$ --- оператор поворота, а $D^{(k)}_{q'q}(\Omega)$ ---
матрица Вигнера ранга $k$. В эквивалентной форме это свойство можно задать
через коммутационные соотношения с компонентами оператора углового момента:

\begin{equation}
[\hat{J}_z,\hat{T}^{(k)}_q]=\hbar q\hat{T}^{(k)}_q,
\end{equation}

\begin{equation}
[\hat{J}_{\pm},\hat{T}^{(k)}_q]=
\hbar\sqrt{k(k+1)-q(q\pm1)}\,\hat{T}^{(k)}_{q\pm1}.
\end{equation}

Эти соотношения показывают, что компоненты тензорного оператора ведут себя
при вращениях так же, как компоненты состояния с угловым моментом $k$. В
частности, скалярный оператор соответствует рангу $k=0$, векторный оператор
--- рангу $k=1$, а квадрупольные и более высокие вклады описываются
тензорами рангов $k=2,3,\ldots$. Для молекулярного гамильтониана особенно
важны скалярные комбинации, поскольку энергия изолированной молекулы не
должна зависеть от ориентации молекулы в лабораторной системе координат. В
то же время оператор дипольного момента является тензором первого ранга, и
именно его симметрийные свойства определяют правила отбора и интенсивности
инфракрасных переходов.

\subsection{Тензорные произведения и матричные элементы}

Основное преимущество неприводимой записи состоит в том, что произведение
двух операторов можно сразу разложить на суммы операторов определенного
ранга. Если $\hat{A}^{(k_1)}$ и $\hat{B}^{(k_2)}$ --- два неприводимых
тензорных оператора, то их связанное произведение записывается как

\begin{equation}
\left[\hat{A}^{(k_1)}\otimes \hat{B}^{(k_2)}\right]^{(k)}_q=
\sum_{q_1q_2}
\langle k_1q_1k_2q_2|kq\rangle
\hat{A}^{(k_1)}_{q_1}\hat{B}^{(k_2)}_{q_2},
\end{equation}

где $\langle k_1q_1k_2q_2|kq\rangle$ --- коэффициенты Клебша--Гордана, а
ранг результирующего оператора удовлетворяет условию
$|k_1-k_2|\leq k\leq k_1+k_2$. При построении эффективного гамильтониана
отбираются прежде всего те комбинации, которые в полной молекулярной группе
дают скалярный вклад. Поэтому тензорная запись естественным образом
связывает аналитический вид оператора с симметрией молекулы и позволяет
заранее исключить запрещенные матричные элементы.

Вычисление матричных элементов тензорных операторов существенно упрощается
благодаря теореме Вигнера--Эккарта. Для базиса, классифицированного по
квантовым числам углового момента, матричный элемент можно представить в
виде произведения чисто геометрического множителя и приведенного матричного
элемента:

\begin{equation}
\langle J'M'|\hat{T}^{(k)}_q|JM\rangle =
(-1)^{J'-M'}
\begin{pmatrix}
J' & k & J\\
-M' & q & M
\end{pmatrix}
\langle J'||\hat{T}^{(k)}||J\rangle .
\end{equation}

Первый множитель содержит только правила сложения угловых моментов и
задает ориентационную часть задачи, а приведенный матричный элемент не
зависит от проекций $M$ и $M'$. В спектроскопии это разделение особенно
важно: оно позволяет отделить универсальную вращательную часть от
молекулярно-специфической информации, связанной с колебательными функциями,
параметрами гамильтониана и параметрами эффективного дипольного момента.

Если вместо полной группы вращений используется конечная точечная группа
молекулы, то аналогичная идея сохраняется. Состояния и операторы
классифицируются по неприводимым представлениям $\Gamma$, а матричный
элемент

\begin{equation}
\langle \Phi_f^{\Gamma_f}|\hat{O}^{\Gamma_O}|\Phi_i^{\Gamma_i}\rangle
\end{equation}

может быть отличен от нуля только в том случае, если прямое произведение
$\Gamma_f\otimes\Gamma_O\otimes\Gamma_i$ содержит полносимметричное
представление. Для гамильтониана таким представлением является $A_1$, а для
оператора дипольного момента --- представление, по которому преобразуются
компоненты вектора дипольного момента в данной точечной группе. Из этого
непосредственно следуют симметрийные правила отбора, используемые при
идентификации спектральных линий.

\subsection{Применение в эффективных гамильтонианах и операторах дипольного момента}

В полуэмпирическом анализе спектров высокого разрешения гамильтониан обычно
не используется в исходном полном виде. Вместо него строится эффективный
гамильтониан, действующий в ограниченном пространстве выбранных
колебательно-вращательных состояний. Его можно представить как сумму
операторных членов различного порядка:

\begin{equation}
\hat{H}_{\mathrm{eff}}=\sum_s h_s\hat{O}_s,
\end{equation}

где $h_s$ --- подгоняемые спектроскопические параметры, а $\hat{O}_s$ ---
операторные структуры, построенные из компонент углового момента,
колебательных координат, импульсов и операторов рождения и уничтожения. В
тензорной форме каждый такой член записывается как скалярная комбинация
неприводимых тензоров. Это позволяет систематически перечислять допустимые
вклады заданного порядка, контролировать их симметрию и избегать избыточных
параметров. Именно поэтому методы эффективных гамильтонианов и операторной
теории возмущений тесно связаны с аппаратом неприводимых тензорных
операторов \cite{Makushkin1984,Papousek1982}.

Аналогичный подход применяется к эффективному оператору дипольного момента,
который определяет интенсивности колебательно-вращательных переходов. В
общем виде лабораторная компонента дипольного момента записывается через
компоненты, связанные с молекулярной системой координат, и направляющие
косинусы:

\begin{equation}
\hat{P}_Z=\sum_{\alpha=x,y,z}k_{Z\alpha}\hat{\mu}_{\alpha}.
\end{equation}

Компоненты $k_{Z\alpha}$ зависят от ориентации молекулы, а
$\hat{\mu}_{\alpha}$ раскладываются по нормальным координатам:

\begin{equation}
\hat{\mu}_{\alpha}=\mu^e_{\alpha}+
\sum_i\mu^i_{\alpha}Q_i+
\frac{1}{2}\sum_{ij}\mu^{ij}_{\alpha}Q_iQ_j+\ldots .
\end{equation}

В этой записи ориентационная часть, колебательная часть и симметрия
оператора разделены. Первый член отвечает за чисто вращательные переходы,
линейные по $Q_i$ члены --- за фундаментальные полосы, квадратичные члены ---
за обертоны и комбинационные полосы. При переходе к эффективному оператору
дипольного момента выполняется унитарное преобразование того же типа, что и
при построении эффективного гамильтониана, поэтому окончательные матричные
элементы интенсивностей выражаются через набор эффективных параметров.
Такая схема используется при определении абсолютных интенсивностей линий и
при переносе информации между изотопологами на основе изотопических
соотношений \cite{Ulenikov2019a,Ulenikov2020a}.

\subsection{Симметризация базисных функций}

Для практического расчета спектра недостаточно только записать оператор в
симметризованном виде. Необходимо также выбрать базисные функции, которые
сами преобразуются по неприводимым представлениям молекулярной группы. В
случае асимметричных и симметричных волчков вращательные функции часто
строятся как линейные комбинации функций $|Jkm\rangle$ с противоположными
значениями $k$. Такие комбинации известны как функции Ванга и позволяют
сразу получить базис с определенной симметрией. В молекулах высокой
симметрии, например в сферических волчках типа $XY_4$ с группой $T_d$,
задача более сложна, поскольку вырожденные колебания требуют согласованного
связывания нескольких угловых моментов и последующего разложения по
неприводимым представлениям точечной группы \cite{hecht}.

Для молекул типа $XY_4$ симметризованную колебательную функцию в
гармоническом приближении удобно записывать как тензорное произведение
частей, относящихся к невырожденному, дважды вырожденному и трижды
вырожденным колебаниям:

\begin{equation}
|v_1;v_2l_2\Gamma_2;v_3l_3,v_4l_4,l n_l\Gamma_{34};n\Gamma\sigma\rangle=
|v_1\rangle
\left(|v_2l_2\Gamma_2\rangle\otimes
|v_3l_3,v_4l_4,l n_l\Gamma_{34}\rangle\right)^{n\Gamma}_{\sigma}.
\end{equation}

Здесь $\Gamma_2$, $\Gamma_{34}$ и $\Gamma$ обозначают симметрию
соответствующих частей функции и полной функции, а $\sigma$ фиксирует строку
неприводимого представления. Часть функции, связанная с двумя трижды
вырожденными колебаниями, может быть представлена через коэффициенты
редукции группы $\mathrm{O}(3)$ к точечной группе молекулы:

\begin{equation}
|v_3l_3,v_4l_4,l n_l\Gamma_{34}\sigma_{34}\rangle=
\sum_{m=-l}^{l}
{}^{l(a)}G^{m}_{n_l\Gamma_{34}\sigma_{34}}
\sum_{|l_3-l_4|\leq l\leq l_3+l_4}
\left(|v_3l_3\rangle\otimes |v_4l_4\rangle\right)^l_m .
\end{equation}

Коэффициенты ${}^{l(a)}G^{m}_{n_l\Gamma_{34}\sigma_{34}}$ выполняют роль
матриц редукции: они переводят базис, естественный для непрерывной группы
вращений, в базис, адаптированный к конечной точечной группе молекулы. В
таком виде симметрия базиса, оператора и матричного элемента задается одной
и той же групповой схемой. Это особенно важно для анализа молекул высокой
симметрии, где одно и то же невозмущенное колебательное состояние может
расщепляться на несколько подуровней различной симметрии, а резонансные
взаимодействия допустимы только между состояниями совместимой симметрии.

\subsection{Значение метода для задач настоящей работы}

В настоящей работе аппарат неприводимых тензорных операторов важен в двух
аспектах. Во-первых, он задает общий язык для построения эффективных
гамильтонианов и операторов дипольного момента, которые используются при
описании положений и интенсивностей линий. Для асимметричных волчков, таких
как изотопологи сероводорода, это проявляется в выборе допустимых
операторных членов, в классификации уровней по симметрии и в записи
матричных элементов, отвечающих за полосы различного типа. При анализе
полосы $\nu_2$ молекул $\mathrm{D_2S}$ и $\mathrm{HDS}$ именно матричные
элементы эффективного дипольного момента позволяют перейти от измеренных
площадей линий к набору параметров, воспроизводящих абсолютные
интенсивности.

Во-вторых, тензорный формализм особенно существенен для молекул высокой
симметрии, рассматриваемых в диссертации. Для тетрагидридов элементов IV
группы симметрия типа $T_d$ приводит к вырожденным колебаниям и к
характерным тетраэдрическим расщеплениям. Без симметризованных
колебательных функций и без тензорного представления операторов построение
матриц гамильтониана становится не только громоздким, но и менее
прозрачным физически. Неприводимая запись позволяет заранее определить,
какие блоки матрицы могут взаимодействовать, какие переходы разрешены
симметрией и какие параметры действительно должны входить в модель. Тем
самым она служит связующим звеном между общей теорией групп, эффективным
гамильтонианом Уотсона, теорией изотопозамещения и практическим анализом
спектров высокого разрешения.

Следовательно, неприводимые тензорные операторы являются не формальным
математическим дополнением, а рабочим инструментом спектроскопического
анализа. Они обеспечивают компактную запись операторов, корректную
симметризацию базиса, систематическое построение матричных элементов и
строгий вывод правил отбора. Благодаря этому данный аппарат используется
как при описании энергетической структуры, так и при анализе интенсивностей
колебательно-вращательных переходов, что делает его необходимой частью
методологической основы настоящей диссертационной работы.

% \section{Методы восстановления потенциальной функции молекулы на основе экспериментальных данных}

% \subsection{Методы ab initio}

% \subsection{Полуэмпирические методы}

